\section{Experimental Design}
\label{sec:experimental-setup}

\subsection{Research Questions}
\researchquestion{RQ1}{What categories of security concerns motivate developer
questions in open-source projects, and which artifact types do maintainers draw
on to answer them?}

\researchquestion{RQ2}{How accurately do frontier LLMs, open-weight LLMs, and
coding agents answer real developer security questions, broken down by question
category and knowledge type?}

\researchquestion{RQ3}{Which artifact types are necessary to answer which
categories of developer security questions? What is the marginal accuracy
contribution of each artifact type when access to it is selectively granted or
withheld?}

\researchquestion{RQ4}{When agents choose autonomously which artifacts to
consult, do they consult the artifact types maintainers actually needed---and
when they fail, is it because they never retrieved the right artifact or
because they reasoned incorrectly over it?}

RQ1 is answered by the benchmark construction itself (\S\ref{sec:open-coding});
the category-to-artifact map it produces is simultaneously an empirical finding
about developer security practice and the ground-truth prediction that RQ3
tests: if maintainers needed a given artifact type to answer questions of a
category, then granting a system access to that artifact type should improve
its accuracy on the category, and withholding it should specifically degrade
it.

\subsection{Evaluation at the Moment of Asking: Time-Capped Snapshots}
\label{sec:snapshot}

A naive with-context evaluation would hand the system the repository at its
current head---but for most items the fix commit, the fix pull-request
discussion, and often the advisory itself entered the repository \emph{after}
the question was asked. Against current state, ``answering'' degenerates into
retrieving the posted resolution. We instead freeze the world at the moment the
question was asked. For an item with report time $t_\theta$, the
\emph{time-capped snapshot} comprises: the repository working tree at the last
default-branch commit preceding $t_\theta$, with commit history truncated
there; the project's issues and pull requests created at or before $t_\theta$,
with later comments and reviews removed; and the GitHub security advisories
referencing the repository published at or before $t_\theta$. The source
thread itself is always excluded from the issue corpus---the system must never
read the conversation it is being tested on. No live network access is
provided under any condition:
a system that consulted today's NVD entry would observe the resolution and void
the cap.

Two properties of this construction deserve emphasis. First, resolution
identifiers that post-date $t_\theta$ are unavailable \emph{by construction},
which changes what can be graded (\S\ref{sec:grading}): for such items the task
is to produce the correct diagnosis, verdict, and remediation direction, not to
name a future release number. Second, when the resolution \emph{pre}-dates the
question---a duplicate report already answered, a fix already merged or
released, an advisory already published---it is legitimately discoverable
inside the snapshot, and locating it is precisely the task a maintainer performs
when triaging a new report. The benchmark therefore contains a natural
retrieval sub-task whose ground truth we did not have to synthesize.

\subsection{Systems Under Test}
\label{sec:systems}

We evaluate three system classes. Each agent backbone is also evaluated as a
bare LLM, so every with-context result has a paired no-context control from the
same model.

\textbf{Bare LLMs (no context).} The model receives only the synthesized
question $q$ and answers from parametric knowledge. We evaluate frontier
closed-weight models and open-weight models under identical prompts
(temperature~0 where the API permits).

\textbf{Reference agent (typed tools).} Our harness implements a tool-calling
agent whose tools are \emph{partitioned by artifact type}: file reading,
listing, and code search over the snapshot working tree (\textsf{code});
commit log and ancestor-guarded commit display (\textsf{commits}); keyword
search and full-text retrieval over the capped issue and pull-request corpora
(\textsf{issues}, \textsf{prs}); and search and retrieval over the capped
advisory set (\textsf{advisory}). Because each tool belongs to exactly one group, the set of
artifact types an agent consulted on an item is read directly off its
transcript as tool-call counts---the instrument RQ4 requires---and access to an
artifact type is granted or withheld by enabling or disabling its tool group,
the instrument RQ3 requires. The agent runs a bounded tool-calling loop
(at most \dnum{15} rounds) before committing to an answer; all transcripts are
released.

\textbf{Off-the-shelf coding agents.} To ground the agent results in systems
practitioners actually use, we additionally evaluate production coding agents
(Claude~Code and OpenCode) headlessly. Each item runs in a sandbox directory
materializing its snapshot as plain files: the repository tree exported at the
snapshot commit \emph{without} version-control metadata (so post-report history
is physically absent), the capped issue, pull-request, and advisory corpora as
structured data files, and a correspondingly truncated commit log. Because a
shell-equipped agent cannot be \emph{provably} prevented from reaching the
network, this condition is best-effort capped: the prompt forbids external
access, the sandbox contains nothing post-report, and we release full
transcripts; we treat the residual risk as a stated limitation
(\S\ref{sec:threats}) and rely on the typed-tool agent for the airtight
analysis.

\subsection{Context Conditions and Selective Artifact Provision}
\label{sec:conditions}

Let $\mathcal{G} = \{\textsf{code}, \textsf{commits}, \textsf{issues},
\textsf{prs}, \textsf{advisory}\}$ be the tool groups. The annotated artifact
types map onto these groups in the obvious way: code, dependency manifests,
and documentation are served by \textsf{code}; advisories and CVE/CWE database
entries by \textsf{advisory}; external references have no serving group, since
live external sources are excluded by the time cap. A \emph{condition} is a
subset $S \subseteq \mathcal{G}$ of groups the system may access:

\begin{itemize}
  \item $S = \emptyset$ --- \textbf{no context}: the bare-LLM condition;
  \item $S = \mathcal{G}$ --- \textbf{full snapshot}: the agent conditions of
        \S\ref{sec:systems}, run over the entire benchmark;
  \item $S = \mathcal{G} \setminus \{g\}$ (leave-one-out) and $S = \{g\}$
        (single-group) --- \textbf{selective provision}, run on a stratified
        subset of contextual items whose annotated artifact types span at least
        two groups, so that ablating one group leaves the others intact.
\end{itemize}

The selective-provision runs operationalize RQ3 as a controlled experiment:
necessity of group $g$ for category $k$ manifests as an accuracy drop under
$\mathcal{G} \setminus \{g\}$ relative to $\mathcal{G}$, and sufficiency as
accuracy under $\{g\}$ approaching that under $\mathcal{G}$
(\S\ref{sec:metrics}). The full-snapshot agent runs simultaneously supply the
RQ4 observational data: per item, the transcript yields the set of groups the
agent actually consulted, which we compare against the groups serving the
maintainer-annotated artifact types.

\subsection{Grading}
\label{sec:grading}

Model answers are free text, so grading combines a deterministic tier that is
exact where exactness is possible with an LLM-judge tier that is validated
where judgment is unavoidable. Both tiers are aware of the condition and the
knowledge type.

\textbf{Tier 1: condition-aware verifiable-fact matching.}
Each verifiable fact is matched in the response by a conservative, type-specific
identifier pattern (word-boundary--guarded version strings, PR-shaped
references, SHA prefixes of at least seven characters, exact CVE/GHSA/CWE
identifiers). Crucially, a fact is \emph{scored} only where it is knowable
under the condition: subject facts the reporter could cite are gradeable under
any condition, while a fix fact is gradeable only when the system had context
access \emph{and} the artifact behind the fact existed at report time (the fix
commit merged, the release tagged, the advisory published before $t_\theta$).
A no-context model is therefore never penalized for not knowing one project's
pull-request numbering, and no system is penalized for not naming an
identifier that did not yet exist---the failure mode that, left uncorrected,
turns a capability benchmark into a memorization probe. Unknowable facts are
recorded as \emph{not gradeable} and excluded from both numerator and
denominator.

\textbf{Tier 2: per-claim judge with a resolution oracle.}
An LLM judge---never the model under evaluation, with candidate identity
anonymized---scores the response against the item's prepared rubric
(\S\ref{sec:rubric-prep}), verdicting each criterion (a cause, a verdict, a
user-actionable remediation, an identifier) as met, partially met, or not met
and quoting the supporting response phrase so every verdict is spot-checkable. Because gold
answers are thread-grounded (\S\ref{sec:normalization}), a terse maintainer
pointer (``fixed in \#932, released 9.0.2'') yields a terse gold; to let a
richer-but-correct response earn credit such a gold cannot confirm, items
whose gold answer cites a fix pull request or commit are enriched at grading
time with the resolved change diff as a \emph{resolution oracle} for the judge.
Eligibility for this oracle is deliberately mechanical---an explicit fix PR or
commit reference, resolved by a single API call, with failures marked
unresolved and the item graded on the thread gold alone---and we report the
eligibility and resolution rates rather than chasing fixes through advisory
prose. The judge is condition-aware in one further respect: under agent
conditions, specific detail beyond the terse gold (affected ranges, root-cause
analysis, identifiers the agent verifiably retrieved) is expected enrichment,
and is flagged as hallucination only when it \emph{contradicts} the gold or the
question's premises; under no context, any concrete project-specific assertion
absent from the gold is flagged.

\textbf{Outcomes.} Each (item, system, condition) triple receives an outcome in
$\{\textsc{correct}, \textsc{partial}, \textsc{incorrect}\}$ by rolling the
per-criterion verdicts up along the rubric axes (\S\ref{sec:rubric-prep}):
correctness \emph{gates} the bucket---an item is \textsc{correct} only when every
correctness criterion is met, \textsc{incorrect} when any is missed or
contradicted, and \textsc{partial} otherwise (correctness only partially met, or
met with thin completeness coverage). Completeness criteria refine within the
bucket and are reported separately as a coverage rate, but never override
correctness, so we need no numeric criterion weights. Orthogonally, a boolean
\textsc{hallucinated} flags a response that asserts fabricated specifics. We keep hallucination separate rather than
folding it into \textsc{incorrect}: an answer can convey the right remediation
\emph{and} fabricate a CVE number, and the fabrication is the more
safety-relevant signal.

\textbf{Judge reliability.} We use an ensemble of $\ge 2$ judge models with
randomized presentation order; disagreements are flagged for human resolution.
Judge verdicts are cross-checked against the Tier-1 identifier matches, and
mismatches (e.g.\ a \textsc{correct} verdict with zero gradeable facts matched)
are audited. We calibrate against human labels on a stratified random sample of
\dnum{50--80} graded items and report Cohen's $\kappa$ (judge vs.\ human); the
judge prompt is revised and re-calibrated until $\kappa \ge 0.80$ before any
reported run.

\subsection{Metrics}
\label{sec:metrics}

Our primary metric is $\mathrm{Acc}(S, k)$, a system's accuracy under
condition $S$ on the items of question category $k$, alongside the analogous
partial-credit rate, the hallucination rate, and the Tier-1 verifiable-fact recall
over gradeable facts. All metrics are additionally split by knowledge type:
no-context capability claims are made on parametric items only, while contextual
items carry the artifact-context analysis. The marginal value of artifact
group $g$ for category $k$ is the leave-one-out contrast on the
selective-provision subset,
$\delta_g(k) = \mathrm{Acc}(\mathcal{G}, k) -
\mathrm{Acc}(\mathcal{G} \setminus \{g\}, k)$, and its sufficiency is
$\sigma_g(k) = \mathrm{Acc}(\{g\}, k) - \mathrm{Acc}(\emptyset, k)$, both over
contextual items. RQ1's category-to-artifact map predicts $\delta_g(k)$ to be
large exactly where maintainers needed an artifact type served by $g$; RQ3
tests that prediction. For RQ4 we report, per category, the agreement between
the groups the agent consulted and the groups serving the maintainer-annotated
artifact types, and decompose agent failures into \emph{retrieval failures}
(the agent answered incorrectly having never consulted some required group)
versus \emph{reasoning failures} (it consulted every required group and still
answered incorrectly).

\subsection{Negative Controls and Contamination Analysis}
\label{sec:contamination}

All threads are public and plausibly present in training corpora; a credible
security benchmark must measure, not assume away, memorization. We use three
instruments. First, the knowledge-type design itself: contextual items are
near-unanswerable without project access, so above-chance no-context accuracy
on contextual items is direct evidence of memorization, cleaner than any
post-hoc heuristic. Second, every item carries $t_\theta$, and results are
stratified by whether $t_\theta$ precedes each model's training cutoff;
markedly higher no-context accuracy on pre-cutoff items indicates recall rather
than reasoning. Third, for the typed-tool agent, the transcript makes
memorization visible at item granularity: a response naming a fix identifier
whose serving tool group was never called cannot have derived it from the
provided context. Approximately \dnum{10\%} of items are additionally tagged as
unanswerable from their annotated artifact types alone (requiring, e.g., an
external database lookup excluded by the time cap); these negative controls
measure confabulation separately from legitimate inference.
